\begin{man}{stdin, stdout, stderr}
    \begin{synopsis}
        \begin{plaincode}[gobble=8]
            #include <stdio.h>

            extern FILE *stdin;
            extern FILE *stdout;
            extern FILE *stderr;
        \end{plaincode}
    \end{synopsis}

    \begin{desc}
        Unter normalen Bedingungen hat jedes UNIX-Betriebssystem drei Kanäle offen, die es erlauben,
        Daten mit einem Benutzer auszutauschen, nämlich
        \begin{itemize}[noitemsep]
            \item den Eingabekanal (z.B. für Tastatureingabe),
            \item den Ausgabekanal (z.B. für Textausgabe auf der Konsole) und
            \item den Fehlerkanal (der nur für die Ausgabe von Fehlertexten bestimmt ist, meistens auch auf der Konsole).
        \end{itemize}
        FIXME
    \end{desc}

    % \begin{example}
    % \end{example}

    % \begin{returns}
    % \end{returns}

    % \begin{seeother}
    % \end{seeother}
\end{man}

